% Options for packages loaded elsewhere
\PassOptionsToPackage{unicode}{hyperref}
\PassOptionsToPackage{hyphens}{url}
\documentclass[
  ignorenonframetext,
]{beamer}
\newif\ifbibliography
\usepackage{pgfpages}
\setbeamertemplate{caption}[numbered]
\setbeamertemplate{caption label separator}{: }
\setbeamercolor{caption name}{fg=normal text.fg}
\beamertemplatenavigationsymbolsempty
% remove section numbering
\setbeamertemplate{part page}{
  \centering
  \begin{beamercolorbox}[sep=16pt,center]{part title}
    \usebeamerfont{part title}\insertpart\par
  \end{beamercolorbox}
}
\setbeamertemplate{section page}{
  \centering
  \begin{beamercolorbox}[sep=12pt,center]{section title}
    \usebeamerfont{section title}\insertsection\par
  \end{beamercolorbox}
}
\setbeamertemplate{subsection page}{
  \centering
  \begin{beamercolorbox}[sep=8pt,center]{subsection title}
    \usebeamerfont{subsection title}\insertsubsection\par
  \end{beamercolorbox}
}
% Prevent slide breaks in the middle of a paragraph
\widowpenalties 1 10000
\raggedbottom
\AtBeginPart{
  \frame{\partpage}
}
\AtBeginSection{
  \ifbibliography
  \else
    \frame{\sectionpage}
  \fi
}
\AtBeginSubsection{
  \frame{\subsectionpage}
}
\usepackage{iftex}
\ifPDFTeX
  \usepackage[T1]{fontenc}
  \usepackage[utf8]{inputenc}
  \usepackage{textcomp} % provide euro and other symbols
\else % if luatex or xetex
  \usepackage{unicode-math} % this also loads fontspec
  \defaultfontfeatures{Scale=MatchLowercase}
  \defaultfontfeatures[\rmfamily]{Ligatures=TeX,Scale=1}
\fi
\usepackage{lmodern}
\ifPDFTeX\else
  % xetex/luatex font selection
\fi
% Use upquote if available, for straight quotes in verbatim environments
\IfFileExists{upquote.sty}{\usepackage{upquote}}{}
\IfFileExists{microtype.sty}{% use microtype if available
  \usepackage[]{microtype}
  \UseMicrotypeSet[protrusion]{basicmath} % disable protrusion for tt fonts
}{}
\makeatletter
\@ifundefined{KOMAClassName}{% if non-KOMA class
  \IfFileExists{parskip.sty}{%
    \usepackage{parskip}
  }{% else
    \setlength{\parindent}{0pt}
    \setlength{\parskip}{6pt plus 2pt minus 1pt}}
}{% if KOMA class
  \KOMAoptions{parskip=half}}
\makeatother
\usepackage{longtable,booktabs,array}
\usepackage{caption}
\captionsetup[table]{skip=6pt}
\newcounter{none} % for unnumbered tables
\usepackage{calc} % for calculating minipage widths
\usepackage{caption}
% Make caption package work with longtable
\makeatletter
\def\fnum@table{\tablename~\thetable}
\makeatother
\setlength{\emergencystretch}{3em} % prevent overfull lines
\providecommand{\tightlist}{%
  \setlength{\itemsep}{0pt}\setlength{\parskip}{0pt}}
% Auto-generated by infrastructure/rendering/_pdf_latex_helpers.py::extract_math_font_preamble.
% Minimal header passed to Pandoc via -H so Beamer slides pick up the same math font
% as the combined-PDF path. Adjust the manuscript's preamble.md to change the font globally.
\usepackage{unicode-math}
\setmathfont{latinmodern-math.otf}

% Slide layout defaults for warning-clean scientific decks.
\usepackage{etoolbox}
\IfFileExists{xurl.sty}{\usepackage{xurl}}{}
\IfFileExists{seqsplit.sty}{\usepackage{seqsplit}}{\newcommand{\seqsplit}[1]{#1}}
\protected\def\breakseq#1{\seqsplit{#1}}
\protected\def\breaktt#1{\begingroup\ttfamily\seqsplit{#1}\endgroup}
\setlength{\emergencystretch}{6em}
\tolerance=5000
\hbadness=10000
\hfuzz=1pt
\setlength{\tabcolsep}{2pt}
\AtBeginEnvironment{longtable}{\tiny\renewcommand{\arraystretch}{0.86}\setlength{\tabcolsep}{1pt}}
\AtBeginEnvironment{tabular}{\tiny\renewcommand{\arraystretch}{0.86}\setlength{\tabcolsep}{1pt}}
\AtBeginEnvironment{equation}{\tiny}
\AtBeginEnvironment{equation*}{\tiny}
\AtBeginEnvironment{align}{\tiny}
\AtBeginEnvironment{align*}{\tiny}
\AtBeginEnvironment{itemize}{\footnotesize}
\AtBeginEnvironment{enumerate}{\footnotesize}
\AtBeginEnvironment{description}{\footnotesize}
\setbeamerfont{caption}{size=\tiny}
\setbeamerfont{caption name}{size=\tiny}
\setbeamerfont{normal text}{size=\small}
\setbeamerfont{frametitle}{size=\small}
\setbeamerfont{section title}{size=\footnotesize}
\setbeamerfont{subsection title}{size=\footnotesize}
\setbeamertemplate{section page}{%
  \centering
  \begin{beamercolorbox}[sep=12pt,center,wd=\paperwidth]{section title}
    \parbox{0.86\paperwidth}{\centering\usebeamerfont{section title}\insertsection\par}
  \end{beamercolorbox}
}
\setbeamertemplate{subsection page}{%
  \centering
  \begin{beamercolorbox}[sep=8pt,center,wd=\paperwidth]{subsection title}
    \parbox{0.86\paperwidth}{\centering\usebeamerfont{subsection title}\insertsubsection\par}
  \end{beamercolorbox}
}
\setlength{\abovecaptionskip}{2pt}
\setlength{\belowcaptionskip}{0pt}

% Natbib and cross-reference fallbacks — slides don't load natbib
% or cleveref, but combined-PDF manuscript prose may emit these
% commands. The fallback renders citations as a bracketed key list
% and cross-references as detokenized labels so slides don't fail on
% undefined control sequences or raw underscores. \providecommand is
% a no-op if packages load later.
\providecommand{\citep}[1]{[#1]}
\providecommand{\citet}[1]{#1}
\providecommand{\citealp}[1]{#1}
\providecommand{\citeauthor}[1]{#1}
\providecommand{\citeyear}[1]{#1}
\providecommand{\cref}[1]{\texttt{\detokenize{#1}}}
\providecommand{\Cref}[1]{\texttt{\detokenize{#1}}}

\newtheorem{proposition}{Proposition}
\newtheorem{hypothesis}{Hypothesis}
\usepackage{bookmark}
\IfFileExists{xurl.sty}{\usepackage{xurl}}{} % add URL line breaks if available
\urlstyle{same}
% fallback for those not using the hyperref driver hyperxmp:
\makeatletter
\@ifundefined{xmpquote}{\newcommand{\xmpquote}[1]{#1}}{}
\makeatother
\hypersetup{
  hidelinks,
  pdfcreator={LaTeX via pandoc}}

\author{\texorpdfstring{}{}}
\date{}

\begin{document}

\section{Conclusion: a recovery-tested bridge with bounded
claims}\label{sec:conclusion}

\begin{frame}[allowframebreaks]{Conclusion: a recovery-tested bridge
with bounded claims}
Active Fedference makes a deliberately narrow bridge between two bodies
of work that are usually discussed with different objects and different
standards of evidence. Active inference describes agents that infer and
act within generative models; federated generalized Bayes describes how
losses, divergences, and variational objectives can make decentralized
inference less sensitive to bad information (Friston et al. 2017; Da
Costa et al. 2020; Bissiri et al. 2016; Knoblauch et al. 2022). This
paper does not dissolve those distinctions. It puts them into one
categorical implementation, identifies the exact point at which they
coincide, and then measures the consequences of moving away from that
point.
\end{frame}

\begin{frame}[allowframebreaks]{The durable result is a recovery
contract}
\protect\phantomsection\label{sec:conclusion-recovery}
The central contribution is the recovery-tested contract. The KL/NLL
client limits of the declared generalized-Bayes construction recover the
closed-form Bayes update, while the zero-robustness server branch
recovers the project log-linear pool, with maximum measured deviations
of 5.55e-17 and 0 (eq.~6;
eq.~7). This is stronger than a verbal analogy
because the limited identities are stated in the formalism, implemented
in the core, and checked by executable invariants. Under the explicit
shared-support, posterior-log-potential, and fixed-weight bridge, the
server pool specializes Eq. 7's message-combination term; it is not an
assertion that the active-inference and robust-Bayes literatures share
all assumptions, objectives, deployment meanings, or the complete source
protocol.

\end{frame}

\begin{frame}[allowframebreaks]{The durable result is a recovery
contract}
That distinction matters historically and methodologically. Logarithmic
pooling has a substantial literature as an aggregation rule for expert
distributions, including its connections to external Bayesianity,
product-of-experts constructions, and KL-based opinion pooling (Genest
and Zidek 1986; Genest et al. 1986; Hinton 2002; Abbas 2009; Dietrich
2021). FedGVI contributes a generalized-Bayes perspective in which the
loss and divergence determine what robustness means (Mildner et al.
2025). The contribution here is therefore a scoped recasting: the
project's categorical log-linear-pool specialization is the non-robust
corner of the implemented family under its stated bridge assumptions,
and the corner becomes a testable boundary condition for future robust
extensions. This does not identify the complete source belief-sharing
protocol with FedGVI.
\end{frame}

\begin{frame}[fragile,allowframebreaks]{What the evidence establishes
away from the corner}
\protect\phantomsection\label{sec:conclusion-evidence}
The standard-Bayes studies provide a necessary baseline rather than
ornamental background. Communication changes mean free energy by 3.3109
nats, Dirichlet learning reduces KL from 3.4231 to 0.0027, and Bayesian
model reduction selects the declared redundant-pruning candidate. These
results recover the declared mechanisms that make belief sharing
scientifically interesting while keeping the source relationship bounded
to the stated categorical protocol (Friston et al. 2024). The extension
studies then show that communication is not automatically beneficial in
every information geometry: disjoint views can make sharing valuable,
whereas a complementary moving-world control can make additional pooling
unnecessary or mildly costly.

\end{frame}

\begin{frame}[fragile,allowframebreaks]{What the evidence establishes
away from the corner}
The contamination study adds a second lesson. The server-side heuristic
is regime-dependent: robust operating points can give up a little
efficiency when contamination is weak and recover that cost when the
declared attack is severe. At the most severe swept rate, the highest
pooled robust mean reaches 0.9880 against the standard pool's 0.6928
(tbl.~3); at the verdict rate, the matched,
BH-adjusted comparison of sec.~7.5 gives 0.9867
against 0.9021. The result is therefore an operating-point contrast, not
a ranking that holds for every contamination rate, attack target, hidden
state, or calibration regime. This interpretation follows the
simulation-study principle that the data-generating mechanism, estimand,
and Monte Carlo unit should be declared before a numerical result is
treated as general evidence (Morris et al. 2019).

\end{frame}

\begin{frame}[fragile,allowframebreaks]{What the evidence establishes
away from the corner}
The three robustness axes give the result its proper meaning. The
client-side generalized-Bayes update is the FedGVI-faithful,
theorem-bearing axis under its matching loss and regularity assumptions.
The server-side \breaktt{robust\_aggregate} is the sharp empirical
heuristic; its proven property here is the recovery limit, not a
transferred client-side bounded-influence theorem. The
\breaktt{variational\_aggregate} is the conservative complement: it
descends a stated aggregation objective and has a proven raw
effective-weight bound, but its objective-backed control does not make
it the peak-accuracy display leader. Separating these axes prevents a
familiar failure in interdisciplinary work: a guarantee proved for one
operator is silently attached to another operator because both are
described with the same word, ``robust.''
\end{frame}

\begin{frame}[allowframebreaks]{Why the bridge matters for active
inference}
\protect\phantomsection\label{sec:conclusion-significance}
Belief sharing is not merely a communication convenience. In an
active-inference system, a shared posterior can change expected
free-energy calculations, model comparison, and subsequent action
selection. A robustification at the fusion step can therefore alter
behavior even when each local generative model is unchanged. Conversely,
a server rule that suppresses an anomalous belief may also suppress a
rare but correct observation. The relevant scientific question is not
simply whether a robust curve rises; it is which assumptions about
competence, independence, support, and action-relevant uncertainty the
fusion rule encodes (Genest and Zidek 1986; Tresp 2000; Heins et al.
2024).

\end{frame}

\begin{frame}[allowframebreaks]{Why the bridge matters for active
inference}
This is why the recovery corner is a useful organizing device. It gives
a common reference behavior before robustness is introduced, makes the
cost of a server intervention measurable, and lets later work compare
new aggregation rules against an interpretable baseline. The bridge also
keeps the distinction between inference and infrastructure visible. The
current federation tests show that serialized beliefs can travel through
the declared local transport and return a bit-identical consensus, but
they do not turn a mathematical aggregation identity into a claim about
secure, fault-tolerant, or privacy- preserving deployment (McMahan et
al. 2017; Blanchard et al. 2017; Pillutla et al. 2022).
\end{frame}

\begin{frame}[allowframebreaks]{What remains unproved}
\protect\phantomsection\label{sec:conclusion-boundaries}
The evidence does not establish universal Byzantine tolerance, truth
recovery, calibration, or an optimal robustness parameter. The primary
intervals are conditional on the fixed hidden state and attack target,
and the nested trial and seed structure is reduced at the declared unit
rather than treated as a larger independent sample (Koehler et al. 2009;
Loy and Korobova 2021). The categorical state space is the object of the
proof and experiment; continuous or hybrid state spaces are a separate
mathematical extension. The neural classification complement uses
point-estimate weights, so it does not establish full posterior FedGVI
behavior at the scale of the source experiments.

\end{frame}

\begin{frame}[allowframebreaks]{What remains unproved}
These are not defects to be hidden by a more expansive title. They
define the proper contribution: an executable categorical bridge, a
recovery certificate, an explicit map of theorem-bearing and heuristic
components, and conditional evidence about contamination behavior.
Robust-statistics language such as influence and breakdown remains
useful for describing the failure modes, but a finite simulation sweep
is not by itself a general estimator-level robustness theorem (Huber and
Ronchetti 2009). The same discipline applies to historical and
conceptual scholarship: early work on probability, inverse inference,
utility, and collective judgment supplies lineage, not evidence for
modern KL, FedGVI, or adversarial-federation claims.
\end{frame}

\begin{frame}[fragile,allowframebreaks]{A falsifiable research program}
\protect\phantomsection\label{sec:conclusion-program}
The next stage should be organized around boundary conditions rather
than a larger collection of demonstrations. First, derive a server
objective whose minimizer is competitive with \breaktt{robust\_aggregate}
while retaining the variational rule's effective-weight control. Recent
work on closed-form generalized variational objectives, logarithmic-pool
weighting, and robust divergence-weighted federation provides relevant
mathematical constraints (Nguyen and Westerhout 2026; Carvalho et al.
2023; Li et al. 2022). A useful candidate must recover the standard
log-linear pool at zero robustness, state which quantity is bounded, and
fail visibly when those assumptions are violated.

\end{frame}

\begin{frame}[fragile,allowframebreaks]{A falsifiable research program}
Second, broaden the primary estimand across attack targets, hidden
states, adaptive adversaries, calibration conditions, and model classes.
The decisive falsifier is not a lower average score on one new grid; it
is failure of the claimed robustness advantage, recovery identity, or
stated uncertainty calibration under a pre-registered extension of the
data-generating mechanism. Third, promote the point-estimate neural
complement to a posterior-parameterized FedGVI study at
source-comparable scale, and test whether the client-side bounded-loss
behavior survives capacity, optimization, and posterior-family changes.
Finally, move from local transport to multi-machine execution with
explicit threat models, authentication, failure handling, and privacy
claims; none should be inferred from the current bit-identity result.
\end{frame}

\begin{frame}[allowframebreaks]{Final position}
\protect\phantomsection\label{sec:conclusion-position}
The strongest conclusion is consequently neither that robust belief
sharing has been solved nor that the bridge is merely metaphorical. In
the declared categorical setting, the standard-Bayes client limits and
project log-linear-pool server identity are tested recovery conditions;
away from them, server behavior is measurable, regime-dependent, and
partitioned into theorem-bearing, objective-backed, and heuristic
claims. That is a modest result, but it is a useful one: it supplies a
reproducible starting point from which stronger theorems, broader state
spaces, independent implementations, and real federated deployments can
be judged without losing the baseline they are meant to extend.
\end{frame}

\end{document}
